Soit $ABC$ un triangle de cercle circonscrit $\Omega$, et soit $O$ le centre de $\Omega$. Un cercle $\Gamma$ de
centre $A$ rencontre le segment $[BC]$ aux points $D$ et $E$, de sorte que $B$, $D$, $E$ et $C$ sont distincts
et dans cet ordre sur la droite $(BC)$. On note $F$ et $G$ les points d'intersection de $\Gamma$ et $\Omega$,
de sorte que $A$, $F$, $B$, $C$ et $G$ sont dans cet ordre sur $\Omega$. Soit $K$ le second point
d'intersection du cercle circonscrit au triangle $BDF$ avec le segment $[AB]$. Soit $L$ le second point
d'intersection du cercle circonscrit au triangle $CGE$ avec le segment $[CA]$.

On suppose que les droites $(FK)$ et $(GL)$ ne sont pas confondues et qu'elles se rencontrent au point $X$.
Prouver que $X$ appartient à la droite $(AO)$.
